% 本文件由 tools/build_book.py 从 dist/reports/assertions.json 生成，请勿手改。
\begin{longtable}{@{}p{0.06\textwidth}p{0.32\textwidth}p{0.22\textwidth}p{0.16\textwidth}@{}}
\caption{章节 6.1 数值断言（生成时间 2026-10-19T16:08:50-0400）}\\
\toprule
类别 & 断言 & 理论值 & 实测值 \\
\midrule
\endfirsthead
\toprule
类别 & 断言 & 理论值 & 实测值 \\
\midrule
\endhead
\bottomrule
\endlastfoot

A & DST Poisson 残差与 Helmholtz 制造解 & max|$\Delta$\_h$\psi$+$\omega$| <=\allowbreak{} 1e-10；max|$\psi$\_H-$\psi$\_exact| <=\allowbreak{} 1e-10 & max|res|=\allowbreak{}4.885e-14，max|$\psi$\_H-$\psi$\_exact|=\allowbreak{}1.110e-15 \checkmark \\
A & 由流函数得到的速度场散度；区间外插值 & 离散散度 <=\allowbreak{} 1e-12(\allowbreak{}五点/\allowbreak{}中心差分同源)；区间外取端点常数 & max|div|=\allowbreak{}8.882e-16；插值 [1.0,\allowbreak{} 1.0,\allowbreak{} 4.0,\allowbreak{} 7.0,\allowbreak{} 7.0] vs [1.0,\allowbreak{} 1.0,\allowbreak{} 4.0,\allowbreak{} 7.0,\allowbreak{} 7.0] \checkmark \\
B & Re=1000/n=33 与 Re=100/n=49 的残差 vs Ghia RMSE & 两组残差都 <1e-6；Re=\allowbreak{}1000 粗网格 RMSE>0.2，而 Re=\allowbreak{}100 细网格 RMSE<1e-2 & Re=\allowbreak{}1000 n=\allowbreak{}33：res=\allowbreak{}9.89e-08，RMSE=\allowbreak{}0.2413；Re=\allowbreak{}100 n=\allowbreak{}49：res=\allowbreak{}9.91e-08，RMSE=\allowbreak{}0.0060 \checkmark \\
A & Re=100 的 Ghia RMSE 网格收敛 & RMSE 单调下降；log(\allowbreak{}RMSE) vs log(\allowbreak{}h) 的斜率在 (\allowbreak{}0.5,\allowbreak{}2.5)(\allowbreak{}一阶迎风约 1) & n=\allowbreak{}17: 0.02024；n=\allowbreak{}25: 0.01211；n=\allowbreak{}33: 0.00870；n=\allowbreak{}49: 0.00603；实测阶 1.10 \checkmark \\
B & Re=50/400 的 dt 限制归属与交叉点 & Re\_cross=\allowbreak{}102.4；Re=\allowbreak{}50 属扩散侧、\allowbreak{}Re=\allowbreak{}400 属 CFL 侧 & Re=\allowbreak{}50: dt\_diff=\allowbreak{}0.01221,\allowbreak{} dt\_cfl=\allowbreak{}0.02500,\allowbreak{} binding=\allowbreak{}0；Re=\allowbreak{}400: dt\_diff=\allowbreak{}0.09766,\allowbreak{} dt\_cfl=\allowbreak{}0.02500,\allowbreak{} binding=\allowbreak{}1 \checkmark \\
B & 中心格式阈值与格点雷诺数 & Re\_th=\allowbreak{}2(\allowbreak{}n-1)；Re=\allowbreak{}1000 在 n=\allowbreak{}33/\allowbreak{}41/\allowbreak{}49 上都已超过阈值；阈值处 Re\_h$\approx$2 & n=\allowbreak{}33: Re\_th=\allowbreak{}64,\allowbreak{} Re\_h(\allowbreak{}Re=\allowbreak{}1000)=\allowbreak{}31.25,\allowbreak{} Re\_h(\allowbreak{}Re\_th)=\allowbreak{}2.000；n=\allowbreak{}41: Re\_th=\allowbreak{}80,\allowbreak{} Re\_h(\allowbreak{}Re=\allowbreak{}1000)=\allowbreak{}25.00,\allowbreak{} Re\_h(\allowbreak{}Re\_th)=\allowbreak{}2.000；n=\allowbreak{}49: Re\_th=\allowbreak{}96,\allowbreak{} Re\_h(\allowbreak{}Re=\allowbreak{}1000)=\allowbreak{}20.83,\allowbreak{} Re\_h(\allowbreak{}Re\_th)=\allowbreak{}2.000 \checkmark \\
A & 四页端点帧的有限性、JSON 常量与 requiredAux & 非有限常量 =\allowbreak{} 0；帧数 =\allowbreak{} 6；每页 requiredAux 都在数组名里 & 6 帧、\allowbreak{}58 个数组，非有限常量计数 0(\allowbreak{}cavity\_vorticity: 10 数组 /\allowbreak{} 33 meta；ghia\_benchmark: 14 数组 /\allowbreak{} 18 meta；turing\_pattern: 7 数组 /\allowbreak{} 34 meta；dam\_break: 10 数组 /\allowbreak{} 20 meta) \checkmark \\
A & 方腔高/低 Re 的残差、RMSE 与时间步绑定 & 高 Re：残差 <=\allowbreak{}1e-6、\allowbreak{}RMSE>0.15、\allowbreak{}CFL 绑定、\allowbreak{}Re\_h>2；低 Re：RMSE<0.02；residual[-1]=\allowbreak{}log10(\allowbreak{}final) & Re=\allowbreak{}1000/\allowbreak{}n=\allowbreak{}41：残差 1.00e-07，RMSE 0.2271，binding=\allowbreak{}1，Re\_h=\allowbreak{}25.0；Re=\allowbreak{}100/\allowbreak{}n=\allowbreak{}49：RMSE 0.0060；residual[-1]=\allowbreak{}-7.000014 vs log10(\allowbreak{}final)=\allowbreak{}-7.000014；中心阈值=\allowbreak{}80.0 \checkmark \\
A & Ghia 四网格 RMSE 曲线与抄写地板 & n\_grid=\allowbreak{}[17,\allowbreak{}25,\allowbreak{}33,\allowbreak{}49]；RMSE 严格递减；斜率$\in$[1.0,\allowbreak{}1.2]；n=\allowbreak{}25 的 RMSE >=\allowbreak{}1.8$\times$n=\allowbreak{}49；中心距 <0.05 & n=\allowbreak{}[17,\allowbreak{} 25,\allowbreak{} 33,\allowbreak{} 49]，RMSE=\allowbreak{}[0.0202 0.0121 0.0087 0.006 ]；斜率=\allowbreak{}1.105，best\_n=\allowbreak{}49；n=\allowbreak{}25 RMSE=\allowbreak{}0.0121 vs n=\allowbreak{}49 0.0060；抄写地板=\allowbreak{}1.0e-02，中心距=\allowbreak{}0.0174；基准点 17 个，u 数组同长=\allowbreak{}True \checkmark \\
A & 四预设的 turing\_\allowbreak{}band / base\_\allowbreak{}valid / $\lambda$(0) / 波长有效性 & turing：band=\allowbreak{}1、\allowbreak{}base=\allowbreak{}1、\allowbreak{}dt\_clamped=\allowbreak{}1、\allowbreak{}波长差<12\%；stripes：band=\allowbreak{}0、\allowbreak{}base=\allowbreak{}1、\allowbreak{}$\lambda$(\allowbreak{}0)>0；spots：band=\allowbreak{}0、\allowbreak{}base=\allowbreak{}0；uniform：band=\allowbreak{}0、\allowbreak{}wavelength\_valid=\allowbreak{}0 且 wavelength=\allowbreak{}0 & turing：$\varphi$*=\allowbreak{}53.899，$\lambda$*=\allowbreak{}0.01502，band=\allowbreak{}1，dt 0.25->0.09537(\allowbreak{}limit 0.09537)，波长 0.1253 vs 0.1166(\allowbreak{}差 7.5\%)；stripes：band=\allowbreak{}0，base=\allowbreak{}1，$\lambda$(\allowbreak{}0)=\allowbreak{}+0.00539；spots：band=\allowbreak{}0，base=\allowbreak{}0；uniform：band=\allowbreak{}0，wavelength\_valid=\allowbreak{}0，wavelength=\allowbreak{}0.000 \checkmark \\
A & 干底溃坝的零步帧、扇区 L1、质量漂移与 CFL & t=\allowbreak{}0：n\_steps=\allowbreak{}0、\allowbreak{}L1=\allowbreak{}0、\allowbreak{}h=\allowbreak{}=\allowbreak{}h\_initial；t=\allowbreak{}0.06：l1\_all>l1\_fan>0、\allowbreak{}0<mass\_drift<1e-6、\allowbreak{}max\_cfl<=\allowbreak{}0.8；n=\allowbreak{}100$\to$400 扇内 L1 递减 & t=\allowbreak{}0：n\_steps=\allowbreak{}0，L1=\allowbreak{}0.0e+00，h 逐位相等=\allowbreak{}True；t=\allowbreak{}0.06：fan=\allowbreak{}113 点，valid=\allowbreak{}166 点，L1\_fan=\allowbreak{}2.0200e-02，L1\_all=\allowbreak{}2.2075e-02，比值=\allowbreak{}1.093；质量漂移=\allowbreak{}9.90e-09，max\_cfl=\allowbreak{}0.800；n=\allowbreak{}100/\allowbreak{}400 扇内 L1=\allowbreak{}3.0410e-02/\allowbreak{}1.2489e-02 \checkmark \\
B & 四页数组与直接调用 numcore 的最大差 & 方腔场/\allowbreak{}Ghia 基准/\allowbreak{}图灵 v/\allowbreak{}溃坝 h 的最大差都 =\allowbreak{} 0 & 方腔 $\omega$ 0.0e+00、\allowbreak{}$\psi$ 0.0e+00、\allowbreak{}x 0.0e+00；Ghia u\_num 0.0e+00、\allowbreak{}u\_ref 0.0e+00；图灵 v 0.0e+00；溃坝 h 0.0e+00、\allowbreak{}x 0.0e+00 \checkmark \\
A & Ghia 三列的表长与端点 & 每列 17 点；u(\allowbreak{}0)=\allowbreak{}0、\allowbreak{}u(\allowbreak{}1)=\allowbreak{}1、\allowbreak{}v(\allowbreak{}0)=\allowbreak{}v(\allowbreak{}1)=\allowbreak{}0；采样严格递增 & Re=\allowbreak{}100: u$\in$[-0.21090,\allowbreak{}1.00000]，v$\in$[-0.24533,\allowbreak{}0.17527]；Re=\allowbreak{}400: u$\in$[-0.32726,\allowbreak{}1.00000]，v$\in$[-0.44993,\allowbreak{}0.30203]；Re=\allowbreak{}1000: u$\in$[-0.38289,\allowbreak{}1.00000]，v$\in$[-0.51550,\allowbreak{}0.37095] \checkmark \\
A & 主涡位置随 Re 的变化与转载分歧地板 & x、\allowbreak{}y 随 Re 单调减小；Re=\allowbreak{}1000 存在两个底角二次涡；分歧地板 1e-2 & Re=\allowbreak{}100 主涡 (\allowbreak{}0.6172,\allowbreak{} 0.7344)，Re=\allowbreak{}400 (\allowbreak{}0.5547,\allowbreak{} 0.6055)，Re=\allowbreak{}1000 (\allowbreak{}0.5313,\allowbreak{} 0.5625)；二次涡 ['bottom\_left',\allowbreak{} 'bottom\_right']；floor=\allowbreak{}0.01 \checkmark \\
A & 图灵底态的两支与稳定性 & u\_-=\allowbreak{}(\allowbreak{}0.25,\allowbreak{}0.3)，det J=\allowbreak{}+0.0045；u\_+ 分支 det<0；spots 无非平凡稳态 & u\_-=\allowbreak{}(\allowbreak{}0.250000,\allowbreak{}0.300000),\allowbreak{} det=\allowbreak{}0.004500；u\_+=\allowbreak{}(\allowbreak{}0.750000,\allowbreak{}0.100000),\allowbreak{} det=\allowbreak{}-0.001500；spots base=\allowbreak{}None \checkmark \\
A & 最不稳定波数与 det 顶点 & $\varphi$*=\allowbreak{}53.899、\allowbreak{}$\lambda$*=\allowbreak{}0.015017；det 顶点 $\varphi$=\allowbreak{}62.75、\allowbreak{}$\lambda$=\allowbreak{}0.013392；经典三预设最快模都在 $\varphi$=\allowbreak{}0，其中 stripes 的 $\lambda$(\allowbreak{}0)=\allowbreak{}+0.0054 是均匀失稳，不是图灵带 & $\varphi$*=\allowbreak{}53.898800，$\lambda$*=\allowbreak{}0.0150170；det 顶点 $\varphi$=\allowbreak{}62.7495，$\lambda$=\allowbreak{}0.0133924，低 10.8\%；经典预设 (\allowbreak{}0.00,\allowbreak{}-0.0350)、\allowbreak{}(\allowbreak{}0.00,\allowbreak{}0.0054)、\allowbreak{}(\allowbreak{}0.00,\allowbreak{}-0.0140)；stripes $\lambda$(\allowbreak{}0)=\allowbreak{}+0.0054 \checkmark \\
B & dt=0.25 在 n=128 的 turing 预设上被夹紧 & limit=\allowbreak{}0.0953674；实际 dt=\allowbreak{}limit，dt\_requested=\allowbreak{}0.25；clamp=\allowbreak{}False 抛 ValueError & limit=\allowbreak{}0.0953674，dt=\allowbreak{}0.0953674，clamped=\allowbreak{}True，requested=\allowbreak{}0.25 \checkmark \\
B & 五步内 $\Sigma$u 变化与 dt$\cdot$$\Sigma$反应 的逐位对账 & 每步差 <=\allowbreak{} 1e-12(\allowbreak{}相对)；扩散项的周期和恒为 0 & predicted=\allowbreak{}761.520552420045，actual=\allowbreak{}761.520552420045，相对差=\allowbreak{}0.00e+00；逐步 ['1.14e-13',\allowbreak{} '1.14e-13',\allowbreak{} '1.14e-13',\allowbreak{} '0.00e+00',\allowbreak{} '0.00e+00'] \checkmark \\
A & 五点模板增长率误差的网格收敛 & 误差按 n\textasciicircum{}-2 下降；n=\allowbreak{}64/\allowbreak{}128/\allowbreak{}256/\allowbreak{}512 的 log-log 斜率约 -2 & 64:9.32e-03、\allowbreak{}128:3.24e-03、\allowbreak{}256:8.68e-04、\allowbreak{}512:2.21e-04；斜率 -1.81 \checkmark \\
B & 经典 spots 的 $\lambda$* 与饱和振幅 & $\lambda$*=\allowbreak{}-0.035<0，但 4000 步后振幅>0.5：自复制瞬态，不是图灵失稳 & $\varphi$*=\allowbreak{}0.000，$\lambda$*=\allowbreak{}-0.0350，振幅=\allowbreak{}0.729，斑点=\allowbreak{}16，t=\allowbreak{}1000.0 \checkmark \\
A & 图灵饱和图案的波长与振幅 & 理论波长 2$\pi$/\allowbreak{}$\varphi$*=\allowbreak{}0.1166；实测应在 12\% 内，振幅>0.2，局部极大>=\allowbreak{}40 & 实测波长=\allowbreak{}0.1253，理论=\allowbreak{}0.1166，相对差=\allowbreak{}7.5\%；振幅=\allowbreak{}0.302，斑点=\allowbreak{}96 \checkmark \\
A & Rusanov 一致性、波速与 t=0 极限 & 相同态通量逐位相等；a=\allowbreak{}|u|+sqrt(\allowbreak{}gh)；t=\allowbreak{}0 时 h=\allowbreak{}1[x<0.5]，u=\allowbreak{}0 & max|$\Delta$F\_h|=\allowbreak{}0.0e+00，max|$\Delta$F\_hu|=\allowbreak{}0.0e+00，波速 [2.9264282071829006,\allowbreak{} 3.681034829318991,\allowbreak{} 4.6794469180700204] \checkmark \\
A & Ritter 干底解的扇区公式与正性 & 扇区公式逐点成立；扇外 h=\allowbreak{}1/\allowbreak{}0；h>=\allowbreak{}0 & 扇内 max|$\Delta$h|=\allowbreak{}0.00e+00，max|$\Delta$u|=\allowbreak{}0.00e+00；h\_min=\allowbreak{}0.0e+00 \checkmark \\
A & 干底溃坝的窗口误差与阶 & 误差单调下降；扇区阶约 0.7-0.9；全域/\allowbreak{}扇区比 >1.3(\allowbreak{}干锋误差) & n=\allowbreak{}100: tight=\allowbreak{}1.91e-02,\allowbreak{} all=\allowbreak{}3.13e-02,\allowbreak{} CFL=\allowbreak{}0.80；n=\allowbreak{}200: tight=\allowbreak{}1.20e-02,\allowbreak{} all=\allowbreak{}2.07e-02,\allowbreak{} CFL=\allowbreak{}0.80；n=\allowbreak{}400: tight=\allowbreak{}7.42e-03,\allowbreak{} all=\allowbreak{}1.30e-02,\allowbreak{} CFL=\allowbreak{}0.80；阶 [0.6691266073450176,\allowbreak{} 0.6942551943035056]；全域/\allowbreak{}扇区=\allowbreak{}1.75 \checkmark \\
B & 周期湿域与干底的质量漂移 & 周期：<=\allowbreak{}1e-12；干底：>1e-10(\allowbreak{}正性地板注入，不是不守恒) & periodic drift=\allowbreak{}2.22e-16；dry drift=\allowbreak{}9.90e-09 \checkmark \\
B & 静水漂移：well-balanced vs 朴素源项 & well-balanced <1e-14；naive >1e-4；两者应差多个数量级 & well-balanced=\allowbreak{}1.89e-16，naive=\allowbreak{}3.30e-04，比值 1.7e+12 \checkmark \\
B & CFL 稳定/失稳的实测边界 & cfl=\allowbreak{}0.8：ht\_end 到达、\allowbreak{}h\_min>0.2；cfl=\allowbreak{}2.0：非有限或 h\_max<1e-6 & stable: reached=\allowbreak{}True,\allowbreak{} h\_min=\allowbreak{}0.426,\allowbreak{} maxCFL=\allowbreak{}0.80；bad: finite=\allowbreak{}False,\allowbreak{} h\_max=\allowbreak{}inf,\allowbreak{} steps=\allowbreak{}830 \checkmark \\
A & 二维径向溃坝的有限性、正性与出流质量 & h\_min>0；出流边界质量/\allowbreak{}能量不增(\allowbreak{}允许舍入 1e-6)；推到 t\_end & t=\allowbreak{}0.0500,\allowbreak{} steps=\allowbreak{}40,\allowbreak{} h\_min=\allowbreak{}1.000e-01,\allowbreak{} mass=\allowbreak{}0.278418 (\allowbreak{}m0=\allowbreak{}0.278418),\allowbreak{} energy=\allowbreak{}0.852330 (\allowbreak{}e0=\allowbreak{}1.011704),\allowbreak{} maxCFL=\allowbreak{}0.30 \checkmark \\
\end{longtable}
